\begin{lemma}[The energy form is a real symmetric bilinear form with quadratic form $\mathcal{E}$]
Let $\mathcal{E}(\cdot)$ be the $[0,+\infty]$-valued energy of \noderef{Energy},
let $\mathcal{E}(\cdot,\cdot)$ be the real-valued form of \noderef{EnergyBil},
and write $v_\alpha$ for the Hermite coefficients of \noderef{HermiteBasis} and
$|\alpha|$ for the multi-index weight of \noderef{MultiIndexSum}.  Then:
\begin{enumerate}
\item for all $v, w \in L^2(\gamma_n)$ with $\mathcal{E}(v) < \infty$ and
$\mathcal{E}(w) < \infty$, the family
$\bigl(|\alpha|^{1/2}\,v_\alpha\,w_\alpha\bigr)_{\alpha \neq 0}$ is summable;
\item for every $v \in L^2(\gamma_n)$ with $\mathcal{E}(v) < \infty$,
$\mathcal{E}(v,v) = \mathcal{E}(v)$, the right-hand side being read as a real
number;
\item for all $v_1, v_2, w \in L^2(\gamma_n)$ with $\mathcal{E}(v_1) < \infty$,
$\mathcal{E}(v_2) < \infty$ and $\mathcal{E}(w) < \infty$, and for every
$c \in \mathbb{R}$,
\[
  \mathcal{E}(v_1+v_2,\,w) = \mathcal{E}(v_1,w) + \mathcal{E}(v_2,w),
  \qquad
  \mathcal{E}(c\,v_1,\,w) = c\,\mathcal{E}(v_1,w),
  \qquad
  \mathcal{E}(v_1,w) = \mathcal{E}(w,v_1);
\]
\item for all $v_1, v_2 \in L^2(\gamma_n)$ with $\mathcal{E}(v_1) < \infty$ and
$\mathcal{E}(v_2) < \infty$, and every $c \in \mathbb{R}$, also
$\mathcal{E}(v_1+v_2) < \infty$ and $\mathcal{E}(c\,v_1) < \infty$.
\end{enumerate}
Clause (2) is the bridge between the $[0,+\infty]$-valued energy, whose
$+\infty$ convention off $\mathcal{V}$ is built into \noderef{Energy}, and the
real-valued form; clauses (3) and (4) say that $\mathcal{E}(\cdot,\cdot)$ is a
symmetric bilinear form on the linear subspace
$\mathcal{V} = \{v : \mathcal{E}(v) < \infty\}$.
\end{lemma}
\begin{proof}
Throughout write, for $\alpha \neq 0$,
\[
  e_\alpha(v) = |\alpha|^{1/2}\,v_\alpha^2 \geq 0 ,
\]
the $\alpha$-th summand of $\mathcal{E}$; the inequality holds because
$|\alpha|$ is a nonnegative integer by \noderef{MultiIndexSum}, so
$|\alpha|^{1/2} \geq 0$, and $v_\alpha^2 \geq 0$.  By \noderef{Energy},
\[
  \mathcal{E}(v) = \sum_{\alpha \neq 0} e_\alpha(v)
  \qquad\text{as an element of } [0,+\infty] ,
\]
and by \noderef{EnergyBil},
$\mathcal{E}(v,w) = \sum_{\alpha\neq0} |\alpha|^{1/2} v_\alpha w_\alpha$.

\emph{Step 1: finiteness of the energy is summability of its summands.}
Since every $e_\alpha(v)$ is nonnegative, the extended-real bridge of the shared
context applies to the family $(e_\alpha(v))_{\alpha \neq 0}$: the sum
$\sum_{\alpha\neq0} e_\alpha(v)$ is finite in $[0,+\infty]$ if and only if the
family is summable as a family of real numbers, and in that case the two
readings of the sum agree.  Hence
\begin{equation*}
  \mathcal{E}(v) < \infty
  \iff \bigl(e_\alpha(v)\bigr)_{\alpha\neq0} \text{ is summable},
  \quad\text{and then}\quad
  \mathcal{E}(v) = \sum_{\alpha\neq0} e_\alpha(v) \in \mathbb{R}. \tag{1}
\end{equation*}

\emph{Step 2: clause (1).}
Let $\mathcal{E}(v), \mathcal{E}(w) < \infty$.  For every $\alpha \neq 0$, the
elementary inequality $|s|\,|t| \leq \tfrac12(s^2+t^2)$ --- which is
$0 \leq (|s|-|t|)^2$ rearranged --- applied with $s = v_\alpha$, $t = w_\alpha$
and multiplied by $|\alpha|^{1/2} \geq 0$ gives
\begin{equation*}
  \bigl||\alpha|^{1/2} v_\alpha w_\alpha\bigr|
    = |\alpha|^{1/2}\,|v_\alpha|\,|w_\alpha|
    \leq \tfrac12\bigl(e_\alpha(v) + e_\alpha(w)\bigr). \tag{2}
\end{equation*}
By Step 1 both $(e_\alpha(v))$ and $(e_\alpha(w))$ are summable, so by linearity
of summation (shared context) the right-hand side of $(2)$ is a summable family
of nonnegative reals; by the comparison test (shared context) the family of
absolute values on the left is then summable, and a family whose absolute values
are summable is itself summable.  This is clause (1).

\emph{Step 3: clause (2).}
Let $\mathcal{E}(v) < \infty$.  Termwise
$|\alpha|^{1/2} v_\alpha v_\alpha = e_\alpha(v)$, so the two series defining
$\mathcal{E}(v,v)$ and $\sum_{\alpha\neq0} e_\alpha(v)$ have identical terms and
hence identical sums.  By the second half of $(1)$ that common value is
$\mathcal{E}(v)$ read as a real number, which is clause (2).

\emph{Step 4: clause (3).}
By \noderef{HermiteBasis} the coefficient map $v \mapsto (v_\alpha)_\alpha$ is
linear, so $(v_1+v_2)_\alpha = (v_1)_\alpha + (v_2)_\alpha$ and
$(c\,v_1)_\alpha = c\,(v_1)_\alpha$ for every $\alpha$.  Therefore, termwise,
\[
  |\alpha|^{1/2}(v_1+v_2)_\alpha w_\alpha
    = |\alpha|^{1/2}(v_1)_\alpha w_\alpha + |\alpha|^{1/2}(v_2)_\alpha w_\alpha ,
  \qquad
  |\alpha|^{1/2}(c\,v_1)_\alpha w_\alpha
    = c\,|\alpha|^{1/2}(v_1)_\alpha w_\alpha .
\]
By Step 2 the two families $(|\alpha|^{1/2}(v_i)_\alpha w_\alpha)_{\alpha\neq0}$,
$i = 1,2$, are summable, so by linearity of summation (shared context) the sum of
the first display is the sum of the two sums, giving additivity, and the sum of
the second display is $c$ times the sum, giving homogeneity.  Symmetry is
immediate termwise from $v_\alpha w_\alpha = w_\alpha v_\alpha$, the two series
again having identical terms.

\emph{Step 5: clause (4).}
Let $\mathcal{E}(v_1), \mathcal{E}(v_2) < \infty$ and $c \in \mathbb{R}$.  From
$0 \leq (s-t)^2$ we get $(s+t)^2 \leq 2s^2 + 2t^2$, so multiplying by
$|\alpha|^{1/2} \geq 0$,
\[
  0 \leq e_\alpha(v_1+v_2)
    = |\alpha|^{1/2}\bigl((v_1)_\alpha + (v_2)_\alpha\bigr)^2
    \leq 2\,e_\alpha(v_1) + 2\,e_\alpha(v_2),
\]
using linearity of the coefficient map as in Step 4.  By Step 1 the right-hand
side is summable, so by the comparison test (shared context)
$(e_\alpha(v_1+v_2))$ is summable, and by the other direction of $(1)$,
$\mathcal{E}(v_1+v_2) < \infty$.  Likewise
$e_\alpha(c\,v_1) = c^2\,e_\alpha(v_1)$, a scalar multiple of a summable family,
hence summable, so $\mathcal{E}(c\,v_1) < \infty$.
\end{proof}
